% Source: https://artofproblemsolving.com/community/c3h1593946_math_kangaroo_2003_level_56_question_30
% Stable identity: aops-topic-1593946-q01
\textbf{At most three colours.} If there were four different colours, one crayon of each would be $4$ crayons with no two of the same colour, contradicting the first condition.

\textbf{At most three crayons of each colour.} If some colour had $4$ crayons, those $4$ together with any fifth crayon would be $5$ crayons with four of the same colour, contradicting the second condition.

So there are at most $3$ colours with at most $3$ crayons each, i.e.\ at most $9$ crayons, with equality only if there are exactly $3$ colours with exactly $3$ crayons each. Since there are $9$ crayons, this is the case, and blue (which occurs) is one of the three colours. Hence there are $\boxed{3}$ blue crayons, answer $\textbf{(B)}$.

(Such a box exists: $3$ blue, $3$ red, $3$ green crayons satisfy both conditions, by the pigeonhole principle for the first.)
